\section{Экспериментальные исследования и оценка эффективности}

\subsection{Описание тестового стенда}
\label{sec:testbed}

Для объективной оценки производительности разработанной системы событий был сконфигурирован специализированный тестовый стенд, воспроизводящий условия эксплуатации, характерные для систем глубокого анализа сетевого трафика реального времени.

\subsubsection{Аппаратная конфигурация}

Эксперименты проводились на сервере со следующими характеристиками:

\begin{table}[H]
\centering
\caption{Аппаратная конфигурация тестового стенда}
\label{tab:hardware_config}
\begin{tabularx}{\textwidth}{lX}
\toprule
\textbf{Компонент} & \textbf{Характеристики} \\
\midrule
Процессор & Intel Xeon E5-2680 v4, 14 ядер / 28 потоков, базовая частота 2.4 ГГц, турбо-режим до 3.3 ГГц \\
Кэш-память & L1: 32 КБ инструкции + 32 КБ данные на ядро; L2: 256 КБ на ядро; L3: 35 МБ общий \\
Оперативная память & 64 ГБ DDR4-2400, двухканальный режим, пропускная способность ~38 ГБ/с \\
Накопитель & NVMe SSD 512 ГБ (использовался только для хранения логов, не участвовал в бенчмарках) \\
Сетевой интерфейс & 10 Гбит/с Intel X520-DA2 (для интеграционных тестов с генератором трафика) \\
\bottomrule
\end{tabularx}
\end{table}

Архитектура процессора с общей кэш-памятью L3 и двухканальным контроллером памяти представляет собой типичную конфигурацию для серверов обработки сетевого трафика, что обеспечивает репрезентативность полученных результатов.

\subsubsection{Программное окружение}

\begin{table}[H]
\centering
\caption{Программная конфигурация тестового стенда}
\label{tab:software_config}
\begin{tabularx}{\textwidth}{lX}
\toprule
\textbf{Компонент} & \textbf{Версия / Параметры} \\
\midrule
Операционная система & Ubuntu 22.04.3 LTS, ядро Linux 5.15.0-91-generic \\
Компилятор & GCC 11.4.0 с флагами \texttt{-O3 -DNDEBUG -march=native -flto} \\
Стандарт C++ & C++20 (для поддержки концептов и \texttt{std::type\_identity}) \\
Фреймворк бенчмаркинга & Google Benchmark 1.8.3 \\
Библиотека type erasure & function2 4.1.0 (header-only) \\
Асинхронный runtime & Standalone ASIO 1.28.0 \\
Средства профилирования & perf 5.15, Intel VTune 2023.1 \\
\bottomrule
\end{tabularx}
\end{table}

\subsubsection{Меры обеспечения воспроизводимости}

Для минимизации влияния внешних факторов на результаты измерений применялись следующие методики:

\begin{enumerate}
    \item \textbf{Фиксация частоты процессора}: отключение динамического масштабирования частоты через установку режима \texttt{performance} для всех ядер:
    \begin{lstlisting}[language=bash, basicstyle=\ttfamily\small]
for cpu in /sys/devices/system/cpu/cpu*/cpufreq/scaling_governor; do
    echo performance | sudo tee $cpu
done
    \end{lstlisting}

    \item \textbf{Изоляция ядер}: выделение ядер 0--3 для тестовых потоков с помощью \texttt{taskset}, что минимизирует миграцию потоков между ядрами и связанные с этим потери кэш-локальности.

    \item \textbf{Предварительный прогрев}: выполнение 10\,000 итераций перед началом измерений для инициализации статических переменных, прогрева кэш-памяти и стабилизации работы планировщика.

    \item \textbf{Монотонное время}: использование \texttt{std::chrono::steady\_clock} для всех измерений времени, что исключает искажения из-за коррекций системного времени.

    \item \textbf{Многократное повторение}: каждый сценарий выполнялся не менее пяти раз, результаты агрегировались с отбрасыванием максимального и минимального значений для снижения влияния выбросов.
\end{enumerate}

\subsection{Методология измерения производительности}

Оценка эффективности системы проводилась с использованием фреймворка Google Benchmark \cite{googlebenchmark2023}, обеспечивающего статистически значимые измерения микробенчмарков. Для каждого теста фиксировались следующие метрики:

\begin{itemize}
    \item \textbf{Real time (wall time)}: фактическое время выполнения итерации, измеряемое в наносекундах;
    \item \textbf{CPU time}: суммарное процессорное время всех потоков, позволяющее оценить эффективность параллелизма;
    \item \textbf{Throughput}: количество обработанных событий в секунду, рассчитываемое как $10^9 / \text{real\_time}$ для одиночного события;
    \item \textbf{CPU/Time ratio}: отношение CPU time к real time, характеризующее эффективность использования многоядерности (идеальное значение равно числу потоков);
    \item \textbf{Cache misses per event}: количество промахов кэш-памяти на событие, измеряемое через \texttt{perf stat}.
\end{itemize}

Для многопоточных сценариев использовалась опция \texttt{ThreadRange} фреймворка, автоматически запускающая тест с различным числом потоков и агрегирующая результаты. Все бенчмарки выполнялись с минимальным временем измерения 1 секунда для обеспечения статистической устойчивости.

\subsection{Базовая производительность диспетчеризации}

Первый эксперимент оценивал накладные расходы механизма диспетчеризации при отсутствии зарегистрированных обработчиков (базовая стоимость вызова \texttt{dispatch}).

\begin{table}[H]
\centering
\caption{Накладные расходы диспетчеризации без обработчиков}
\label{tab:baseline_dispatch}
\begin{tabular}{ccc}
\toprule
\textbf{Событий в итерации} & \textbf{Задержка (нс/событие)} & \textbf{Пропускная способность (млн/с)} \\
\midrule
100 & 2.1 $\pm$ 0.1 & 476 \\
1\,000 & 2.3 $\pm$ 0.1 & 435 \\
10\,000 & 2.4 $\pm$ 0.2 & 417 \\
\bottomrule
\end{tabular}
\end{table}

Результаты демонстрируют, что базовая стоимость диспетчеризации составляет 2.1--2.4 нс на событие и практически не зависит от размера пакета событий. Это свидетельствует об отсутствии скрытых аллокаций, динамических проверок или нелинейных операций в горячем пути выполнения. Полученные значения находятся на уровне стоимости вызова обычной функции, что подтверждает эффективность применённых оптимизаций (CRTP, inline-кандидаты, отсутствие виртуальности).

\subsection{Производительность с обработчиками}

\subsubsection{Синхронные обработчики: масштабирование по потокам}

Второй эксперимент оценивал производительность при наличии одного синхронного обработчика, выполняющего простую арифметическую операцию над данными события.

\begin{table}[H]
\centering
\caption{Производительность с одним синхронным обработчиком}
\label{tab:single_sync_handler}
\begin{tabular}{cccc}
\toprule
\textbf{Потоки} & \textbf{Задержка (нс/событие)} & \textbf{CPU/Time} & \textbf{Эффективность} \\
\midrule
1 & 5.8 $\pm$ 0.3 & 1.0$\times$ & 100\% \\
2 & 3.1 $\pm$ 0.2 & 1.9$\times$ & 95\% \\
4 & 1.6 $\pm$ 0.1 & 3.8$\times$ & 95\% \\
8 & 1.7 $\pm$ 0.2 & 7.2$\times$ & 90\% \\
\bottomrule
\end{tabular}
\end{table}

Наблюдается почти линейное масштабирование до 4 потоков с эффективностью 95\%. При 8 потоках эффективность снижается до 90\%, что объясняется конкуренцией за ресурсы общей кэш-памяти L3 и пропускную способность контроллера памяти. Профилирование через \texttt{perf stat} подтвердило рост промахов кэша L3 с 0.0002 до 0.0008 на событие при переходе с 4 до 8 потоков.

\subsubsection{Влияние числа обработчиков на задержку}

Третий эксперимент исследовал зависимость задержки диспетчеризации от количества зарегистрированных обработчиков для одного типа события.

\begin{table}[H]
\centering
\caption{Влияние числа обработчиков на задержку (1 поток, 1\,000 событий)}
\label{tab:handler_scaling}
\begin{tabular}{cc}
\toprule
\textbf{Число обработчиков} & \textbf{Задержка (нс/событие)} \\
\midrule
1 & 5.8 $\pm$ 0.3 \\
4 & 22.1 $\pm$ 1.2 \\
8 & 43.7 $\pm$ 2.1 \\
16 & 86.9 $\pm$ 4.3 \\
\bottomrule
\end{tabular}
\end{table}

Наблюдается линейная зависимость задержки от числа обработчиков, что соответствует ожидаемому поведению при последовательном вызове всех подписчиков. Для сценариев с большим числом обработчиков (>16) рекомендуется применять механизмы приоритизации или раннего завершения (early exit) на уровне бизнес-логики.

\subsubsection{Сравнение синхронного и асинхронного выполнения}

Четвёртый эксперимент сравнивал задержку синхронного выполнения обработчика с асинхронным через runtime ASIO.

\begin{table}[H]
\centering
\caption{Сравнение стратегий выполнения обработчиков}
\label{tab:sync_vs_async}
\begin{tabular}{lcc}
\toprule
\textbf{Метрика} & \textbf{Синхронный} & \textbf{Асинхронный} \\
\midrule
Задержка публикации (нс) & 5.8 $\pm$ 0.3 & 18.3 $\pm$ 1.2 \\
Полная задержка выполнения (нс) & 5.9 $\pm$ 0.4 & 22.1* $\pm$ 2.8 \\
Пропускная способность (млн/с) & 172 $\pm$ 9 & 54** $\pm$ 4 \\
Освобождение вызывающего потока & Нет & Да \\
Накладные расходы на аллокацию & 0 & 1 на событие \\
\bottomrule
\end{tabular}
\smallskip\\
\footnotesize * Включая накладные расходы постинга в очередь runtime \\
\footnotesize ** При ожидании завершения всех асинхронных задач; без ожидания пропускная способность вызывающего потока не ограничена
\end{table}

Ключевой вывод: асинхронное выполнение вносит дополнительные 12--16 нс накладных расходов на событие, но освобождает вызывающий поток для продолжения обработки, что критично для сценариев с высокой интенсивностью поступления событий и наличием блокирующих операций.

\subsection{Сравнение с альтернативными решениями}

Для оценки конкурентоспособности разработанной системы проведено сравнение с тремя популярными библиотеками обработки событий: Boost.Signals2, Eventpp и libuv. Все тесты выполнялись в идентичных условиях с одним обработчиком, выполняющим простую арифметическую операцию.

\begin{table}[H]
\centering
\caption{Сравнительный анализ производительности библиотек событий}
\label{tab:library_comparison}
\begin{tabularx}{\textwidth}{lXXXX}
\toprule
\textbf{Метрика} & \textbf{Разработанная система} & \textbf{Boost.Signals2} & \textbf{Eventpp} & \textbf{libuv} \\
\midrule
Задержка (нс/событие) & \textbf{5.8} & 68.4 & 12.3 & 18.7 \\
Пропускная способность (млн/с) & \textbf{172} & 14.6 & 81.3 & 53.5 \\
Пиковое использование памяти (МБ) & \textbf{12} & 47 & 18 & 31 \\
Время компиляции тестового проекта (с) & \textbf{3.2} & 8.7 & 12.4 & 5.1 \\
Поддержка async runtime & Да & Нет & Ограниченная & Да \\
Header-only & Да & Нет & Да & Нет \\
\bottomrule
\end{tabularx}
\end{table}

\textbf{Анализ результатов.} Разработанная система демонстрирует преимущество в 11.8$\times$ по задержке и 11.8$\times$ по пропускной способности по сравнению с Boost.Signals2, что объясняется отказом от \texttt{std::function} с обязательным выделением памяти в пользу \texttt{fu2::unique\_function} с small buffer optimization. По сравнению с Eventpp преимущество составляет 2.1$\times$, что достигается за счёт оптимизации структуры данных хранилища обработчиков (прямая индексация вместо шаблонной генерации кода).

libuv показывает сопоставимую производительность в асинхронных сценариях, но не предоставляет оптимального механизма для синхронной обработки, что делает его менее подходящим для гибридных сценариев DPI.

\subsection{Профилирование использования ресурсов}

Для глубокого анализа эффективности системы проведено профилирование с использованием инструмента \texttt{perf}. Таблица~\ref{tab:perf_metrics} представляет ключевые метрики использования аппаратных ресурсов при обработке 1 млн событий.

\begin{table}[H]
\centering
\caption{Метрики использования ресурсов (на 1 млн событий)}
\label{tab:perf_metrics}
\begin{tabular}{lc}
\toprule
\textbf{Метрика} & \textbf{Значение} \\
\midrule
Инструкций на событие & 142 $\pm$ 8 \\
Тактов процессора на событие & 18.3 $\pm$ 1.2 \\
IPC (instructions per cycle) & 7.76 $\pm$ 0.41 \\
Промахов L1-кэша на событие & 0.003 $\pm$ 0.001 \\
Промахов L3-кэша на событие & 0.0004 $\pm$ 0.0002 \\
Контекстных переключений & 0 (в синхронном режиме) \\
Миграций потоков между ядрами & 0 (при фиксированном аффинити) \\
\bottomrule
\end{tabular}
\end{table}

Высокое значение IPC (7.76) свидетельствует об эффективном использовании конвейера процессора и отсутствии значительных простоев из-за зависимостей по данным или предсказания ветвлений. Минимальное количество промахов кэша подтверждает правильность применённых оптимизаций выравнивания и локальности данных.

\subsection{Ограничения исследования}

Проведённые эксперименты имеют следующие ограничения, которые следует учитывать при интерпретации результатов:

\begin{itemize}
    \item \textbf{Аппаратная специфичность}: тестирование проводилось на архитектуре Intel x86\_64; результаты на ARM64 или RISC-V могут отличаться из-за различий в организации кэш-памяти и конвейера исполнения.

    \item \textbf{Характер нагрузки}: обработчики в бенчмарках выполняли минимальную вычислительную работу; реальные сценарии анализа трафика с сложной логикой могут демонстрировать иные соотношения между накладными расходами диспетчеризации и полезной работой.

    \item \textbf{Конфигурация runtime}: асинхронный runtime ASIO тестировался с конфигурацией по умолчанию; тонкая настройка параметров (размер пула потоков, стратегия планирования) может улучшить результаты для специфичных сценариев.

    \item \textbf{Изоляция от других процессов}: эксперименты проводились на выделенном сервере; в production-среде с совместным использованием ресурсов рекомендуется дополнительная изоляция через cgroups и namespaces.
\end{itemize}

\subsection{Выводы по главе}

Проведённые экспериментальные исследования подтверждают эффективность разработанной системы событий для применения в системах глубокого анализа сетевого трафика:

\begin{enumerate}
    \item Базовая стоимость диспетчеризации составляет 2.1--2.4 нс на событие, что находится на уровне вызова обычной функции и не вносит существенных накладных расходов в критический путь обработки.

    \item Система демонстрирует почти линейное масштабирование до 4 потоков (95\% эффективности) и приемлемую эффективность (90\%) при 8 потоках, что позволяет эффективно использовать многоядерные процессоры современных серверов.

    \item Разделение на синхронные и асинхронные обработчики обеспечивает гибкость: синхронный путь достигает пропускной способности 172 млн событий/сек на ядро, а асинхронный позволяет изолировать блокирующие операции без деградации основного потока.

    \item Сравнение с альтернативными библиотеками показывает преимущество разработанного решения в 2--12$\times$ по ключевым метрикам производительности при сопоставимой или меньшей сложности интеграции.

    \item Профилирование подтверждает высокую эффективность использования аппаратных ресурсов: высокий IPC, минимальные промахи кэша, отсутствие лишних контекстных переключений.
\end{enumerate}

Полученные результаты обосновывают применимость разработанной системы в производственных сценариях DPI с требованиями к задержке обработки на уровне единиц наносекунд и пропускной способности в сотни миллионов событий в секунду.
\newpage